\section{Einführung}
\label{chap:introduction}

Ein portabler Werkzeugordner ist kein Produkt-Template. Er darf vorhandene Anwendung,
Domänendaten und unbekannte Dateien nicht zu seinem Eigentum erklären. Die vorliegende
Fallstudie beantwortet daher die Frage, wie ein kopierbares Tooling mehrere Projektformen
unterstützen kann, ohne Produktquellen zu besitzen und ohne Änderungen unsichtbar zu
machen. Die Methodik folgt einer nachvollziehbaren Fallstudienperspektive
\parencite{runeson2009case-study}.

Untersucht werden Kontextauflösung, Profile, Adapter, Transaktionen, State, Migration,
Export und Tests der aktuellen Tooling-Revision \parencite{template-tooling-source}.
Die folgenden Kapitel unterscheiden beobachtete Implementierung, automatisierte
Verifikation, künftige Arbeit und Grenzen ausdrücklich.
